\subsubsection{Changes between inputs in the Claude batch}
All fifteen judgments that changed between inputs in the Claude batch concern the software and documentation criteria. For TruthfulQA and DesignQA in both repetitions, A1 and A2 move from 5 for Input 1 to 10 for Input 2. For Input 1 those responses state that the implementation cannot be inspected and apply anchor 5; for Input 2 they describe the entry point or compare the implementation with the documented rule and apply anchor 10. For HumanEval, where the paper itself prints the estimator code, A2 is 10 in every session, and A1 is 5 in the first Input 1 repetition and 10 in the second, so the same README produced both anchors within one configuration. No A1 or A2 judgment reaches 15, and each Input 2 response explains why execution for the current edition is not documented. In the Codex batch the Input 1 sessions had used 10 for A1 from a README command alone in four of six cases. The Claude judgments follow the anchor distinction between a claimed and an inspectable implementation more often, but this is a statement about adherence to anchor wording, not about the correctness of either rating.

The remaining changes concern A3, A4 and A5 and are not uniform. For DesignQA A3, the Input 2 response reads CSV headers and confirms metadata fields that the Input 1 response could not, which is a change supported by information only available in the repository. For TruthfulQA A3 the direction differs between repetitions, which indicates a sufficiency threshold about edition metadata applied to text present in both inputs. No measurement criterion changed with access in this batch. Repository files in Input 2 address what can be inspected; they did not alter the responses' reading of construct, checking, reliability, human comparison or exposure evidence.
